\setcounter{section}{11}

  \zwischenueberschrift{Produkte von Mannigfaltigkeiten}

\inputdefinition
{}
{

Es seien
\mathkor {} {M} {und} {N} {}
zwei
\definitionsverweis {differenzierbare Mannigfaltigkeiten}{}{}
mit den
\definitionsverweis {Atlanten}{}{}
\mathkor {} {(U_i,U_i',\alpha_i, i \in I)} {und} {(V_j,V_j',\beta_j, j \in J)} {.}
Dann nennt man den
\definitionsverweis {Produktraum}{}{}

\mathl{M \times N}{} mit den
\definitionsverweis {Karten}{}{}
\maabbdisp {\alpha_i \times \beta_j} {U_i \times V_j } { U_i' \times V_j'
} {}
\zusatzklammer {mit

\mavergleichskettek
{\vergleichskettek
{  (i,j)
}
{ \in }{ I \times J
}
{  }{
}
{    }{
}
{   }{
}
}
{}{}{}
und

\mavergleichskettek
{\vergleichskettek
{ U_i' \times V_j'
}
{ \subseteq }{ \R^m \times \R^n
}
{  }{
}
{    }{
}
{   }{
}
}
{}{}{}} {} {}
das \definitionswort {Produkt der Mannigfaltigkeiten}{}
\mathkor {} {M} {und} {N} {.}

}

Es handelt sich dabei in der Tat um eine Mannigfaltigkeit, siehe
[[Produkt von Mannigfaltigkeiten/Ist Mannigfaltigkeit/Aufgabe|Aufgabe 11.1]].
Eine Produktmannigfaltigkeit der Form
\mathl{M \times \R}{} nennt man auch \stichwort {Zylinder} {} über $M$. Zu abgeschlossenen Untermannigfaltigkeiten

\mavergleichskette
{\vergleichskette
{ M
}
{ \subseteq }{ \R^\ell
}
{  }{
}
{    }{
}
{   }{
}
}
{}{}{}
und

\mavergleichskette
{\vergleichskette
{ N
}
{ \subseteq }{ \R^k
}
{  }{
}
{    }{
}
{   }{
}
}
{}{}{}
ist die Produktmannigfaltigkeit
\mathl{M \times N}{} eine abgeschlossene Untermannigfaltigkeit von

\mavergleichskette
{\vergleichskette
{ \R^\ell \times \R^k
}
{ = }{ \R^{\ell + k}
}
{  }{
}
{    }{
}
{   }{
}
}
{}{}{,}
siehe
[[Zwei abgeschlossene Untermannigfaltigkeiten/Produkt davon/Aufgabe|Aufgabe 11.2]].

__NOEDITSECTION__

\inputfaktbeweis
{Produkt von Mannigfaltigkeiten/Abbildungseigenschaften/Fakt}
{Lemma}
{}
{

\faktsituation {Es seien
\mathkor {} {M} {und} {N} {}
\definitionsverweis {differenzierbare Mannigfaltigkeiten}{}{}
und
\mathl{M \times N}{} ihr
\definitionsverweis {Produkt}{}{.}}
\faktuebergang {Dann gelten folgende Eigenschaften.}
\faktfolgerung {\aufzaehlungdrei{Die
\definitionsverweis {Projektionen}{}{}
\maabbeledisp {p_1} {M \times N} {M
} {(x,y)} { x
} {,}
und
\maabbeledisp {p_2} {M \times N} {N
} {(x,y)} { y
} {,}
sind
\definitionsverweis {differenzierbare Abbildungen}{}{.}
}{Der
\definitionsverweis {Tangentialraum}{}{}
in einem Punkt

\mavergleichskette
{\vergleichskette
{ R
}
{ = }{ (P,Q)
}
{  }{
}
{    }{
}
{   }{
}
}
{}{}{}
ist

\mavergleichskette
{\vergleichskette
{ T_R(M \times N)
}
{ = }{ T_PM \times T_QN
}
{  }{
}
{    }{
}
{   }{
}
}
{}{}{.}
}{Es sei $L$ eine weitere differenzierbare Mannigfaltigkeit. Dann ist eine Abbildung
\maabbeledisp {\varphi \times \psi} { L } { M \times N
} { u } { (\varphi(u), \psi(u))
} {,}
genau dann differenzierbar, wenn die
\definitionsverweis {Komponentenabbildungen}{}{}
\mathkor {} {\varphi} {und} {\psi} {}
differenzierbar sind.
}}
\faktzusatz {}
\faktzusatz {}

}
{

\teilbeweis {}{}{}
{(1). Durch Übergang zu Karten können wir annehmen, dass
\mathkor {} {M} {und} {N} {}
\definitionsverweis {offene Teilmengen}{}{}
im
\mathkor {} {\R^m} {bzw. im} {\R^n} {}
sind. In diesem Fall handelt es sich um eine
\definitionsverweis {Einschränkung}{}{}
der
\definitionsverweis {linearen Projektion}{}{}
\maabb {} {\R^m \times \R^n } {\R^m
} {,}
die
nach [[Differenzierbarkeit/K/Lineare Abbildungen sind differenzierbar/Fakt|Proposition 45.3 (Analysis (Osnabrück 2021-2023))]]
\definitionsverweis {stetig differenzierbar}{}{}
ist.}
{}
\teilbeweis {}{}{}
{(2). Die differenzierbaren Projektionen
\maabb {p_1} {M \times N} {M
} {}
und
\maabb {p_2} {M \times N} {N
} {}
liefern die linearen
\definitionsverweis {Tangentialabbildungen}{}{}
\maabb {T_R(p_1)} {T_R(M \times N)} {T_PM
} {}
und
\maabb {T_R(p_2)} {T_R(M \times N)} {T_QN
} {}
und  damit insgesamt die lineare Abbildung
\maabbdisp {T_R(p_1) \times T_R(p_2)} { T_R(M \times N)} {T_PM \times T_QN} {.}
Zum Nachweis der Bijektivität kann man zu Karten übergehen und annehmen, dass

\mavergleichskette
{\vergleichskette
{ M
}
{ \subseteq }{ \R^m
}
{  }{
}
{    }{
}
{   }{
}
}
{}{}{}
und

\mavergleichskette
{\vergleichskette
{ N
}
{ \subseteq }{ \R^n
}
{  }{
}
{    }{
}
{   }{
}
}
{}{}{}
offene Teilmengen sind. Diese Abbildung wird dann zur Bijektion
\maabbdisp {p_1 \times p_2} {\R^{m+n} } { \R^m \times \R^n
} {.}}
{}
\teilbeweis {}{}{}
{(3). Für einen fixierten Punkt

\mavergleichskette
{\vergleichskette
{ A
}
{ \in }{ L
}
{  }{
}
{    }{
}
{   }{
}
}
{}{}{}
kann man unter Verwendung von Kartenumgebungen von $A$ und von
\mathkor {} {\varphi(A)} {und} {\psi(A)} {}
sich darauf zurückziehen, dass alle drei Mannigfaltigkeiten offene Mengen in euklidischen Räumen sind. Wenn beide Abbildungen stetig differenzierbar sind, so folgt nach
[[Totales Differential/K/Produktabbildungen/Aufgabe|Aufgabe 45.9 (Analysis (Osnabrück 2021-2023))]]
die stetige Differenzierbarkeit der Gesamtabbildung. Die Umkehrung ist klar.}
{}

}

\bild{ \begin{center}
\includegraphics[width=5.5cm]{ \bildeinlesungpng {Toroidal_coord} {png} }
\end{center}
\bildtext {} }

\bildlizenz { Toroidal coord.png } {} {Dave Burke} {Commons} {CC-by-sa 3.0} {}

\inputbeispiel{}
{

Das
\definitionsverweis {Produkt}{}{}
der
\definitionsverweis {Kreislinie}{}{}
mit sich selbst, also

\mavergleichskette
{\vergleichskette
{ M
}
{ = }{ S^1 \times S^1
}
{  }{
}
{    }{
}
{   }{
}
}
{}{}{,}
heißt \stichwort {Torus} {.} Dies ist eine zweidimensionale Mannigfaltigkeit. Da

\mavergleichskette
{\vergleichskette
{ S^1
}
{ \subset }{ \R^2
}
{  }{
}
{    }{
}
{   }{
}
}
{}{}{}
eine
\definitionsverweis {abgeschlossene Untermannigfaltigkeit}{}{}
ist, lässt sich der Torus als abgeschlossenene Untermannigfaltigkeit im

\mavergleichskette
{\vergleichskette
{ \R^2 \times \R^2
}
{ = }{ \R^4
}
{  }{
}
{    }{
}
{   }{
}
}
{}{}{}
realisieren. Sie lässt sich aber auch als abgeschlossenene Untermannigfaltigkeit im $\R^3$ realisieren. Dazu seien
\mathkor {} {r} {und} {R} {}
positive reelle Zahlen mit

\mavergleichskette
{\vergleichskette
{ 0
}
{ < }{ r
}
{ < }{ R
}
{    }{
}
{   }{
}
}
{}{}{.}
Dann ist die Menge

\mathdisp {\left\{ (x,y,z) \in \R^3  \mid   \left(  \sqrt{x^2+y^2} -R  \right)^2 +z^2 = r^2         \right\}} {  }

ein Torus. Es handelt sich bei dieser Realisierung um die Oberfläche eines
\zusatzklammer {aufgeblasenen} {} {}
\anfuehrung{Fahrradschlauches}{,} dessen \anfuehrung{Radradius}{} gleich $R$ und dessen \anfuehrung{Schlauchradius}{} gleich $r$ ist
\zusatzklammer {das Rad liegt in der \mathlk{x-y}{-}Ebene} {} {.}
Der Zusammenhang mit dem Produkt
\mathl{S^1 \times S^1}{} ergibt sich, indem man dem Produktwinkel
\mathl{( \varphi, \psi)}{} den Punkt
\mathl{( (R +r \cos  \psi) \cos \varphi, (R+ r \cos  \psi ) \sin \varphi ,  r \sin  \psi )}{} zuordnet.

}

  \zwischenueberschrift{Die Lie-Klammer von Vektorfeldern}

\inputdefinition
{}
{

Es sei $M$ eine
\definitionsverweis {differenzierbare Mannigfaltigkeit}{}{.}
Eine Abbildung
\maabbdisp {D} {C^1(M,\R) } { C^0(M,\R)
} {}
heißt
\definitionswort {Differentialoperator erster Ordnung}{,}
wenn sie folgende Eigenschaften erfüllt.
\aufzaehlungzwei { $D$ ist
$\R$-\definitionsverweis {linear}{}{.}
} {Es ist

\mavergleichskette
{\vergleichskette
{ D(fg)
}
{ = }{ fD(g)+gD(f)
}
{  }{
}
{    }{
}
{   }{
}
}
{}{}{.}
}

}

Statt Differentialoperator erster Ordnung sagt man auch \stichwort {Derivation} {.}

Zu einer offenen Menge

\mavergleichskette
{\vergleichskette
{ U
}
{ \subseteq }{ \R^n
}
{  }{
}
{    }{
}
{   }{
}
}
{}{}{}
sind die partiellen Ableitungen $\partial_i$ und $\partial_j$ als Operationen auf zweifach stetig differenzierbaren Funktionen vertauschbar, also

\mavergleichskettedisp
{\vergleichskette
{ \partial_i \circ \partial_j
}
{ =} { \partial_j \circ \partial_i
}
{ } {
}
{ } {
}
{ } {
}
}
{}{}{.}
Eine Linearkombination

\mavergleichskettedisp
{\vergleichskette
{ V
}
{ =} { \sum_{i = 1}^n f_i \partial_i
}
{ } {
}
{ } {
}
{ } {
}
}
{}{}{}
mit auf $U$ definierten reellwertigen Funktionen $f_i$ kann man als
\definitionsverweis {Vektorfeld}{}{}
\maabbeledisp {} { U } { TU = U \times \R^n
} { P } { (P, \sum_{i = 1}^n f_i(P) \partial_i ) = (P,f_1(P) , \ldots , f_n(P) )
} {,}
und ebenso als Ableitungsoperator
\maabbeledisp {} {C^1(U,\R)} {C^0(U,\R)
} {h } { \sum_{i = 1}^n f_i  \partial_i  h
} {,}
auffassen
\zusatzklammer {oder von $C^\infty(U,\R)$ nach $C^\infty(U,\R)$} {} {.}
In diesem Sinne entsprechen sich Vektorfelder und Differentialoperatoren der ersten Ordnung. Differentialoperatoren kann man aber darüber hinaus miteinander verknüpfen, wobei sich herausstellt, dass es auf die Reihenfolge ankommt.

__NOEDITSECTION__

\inputfaktbeweis
{Offene Menge/R^n/Vektorfelder/Hintereinanderschaltung/Explizit/Fakt}
{Lemma}
{}
{

\faktsituation {Es sei

\mavergleichskette
{\vergleichskette
{ U
}
{ \subseteq }{ \R^n
}
{  }{
}
{    }{
}
{   }{
}
}
{}{}{}
\definitionsverweis {offen}{}{}
und seien $f_1 , \ldots , f_n, g_1 , \ldots , g_n,h$ zweifach
\definitionsverweis {stetig differenzierbare Funktionen}{}{}
auf $U$.}
\faktfolgerung {Dann gilt

\mavergleichskettedisphandlinks
{\vergleichskettedisphandlinks
{ \left(  \sum_{i = 1}^n f_i \partial_i  \right) \left(  \left(  \sum_{j = 1}^n g_j \partial_j   \right) (h)  \right)
}
{ =} { \sum_{j = 1}^n  \left(  \sum_{i = 1}^n  f_i \partial_i g_j  \right) \partial_j h + \sum_{j = 1}^n \sum_{i = 1}^n    f_i g_j \partial_i \partial_j h
}
{ } {
}
{ } {
}
{ } {
}
}
{}{}{.}}
\faktzusatz {}
\faktzusatz {}

}
{

Es ist

\mavergleichskettealignhandlinks
{\vergleichskettealignhandlinks
{ \left(  \sum_{i = 1}^n f_i \partial_i  \right) \left(  \left(  \sum_{j = 1}^n g_j \partial_j   \right) (h)  \right)
}
{ =} { \left(  \sum_{i = 1}^n f_i \partial_i  \right) \left(  \sum_{j = 1}^n g_j \partial_j (h)  \right)
}
{ =} { \sum_{i = 1}^n f_i \partial_i \left(  \sum_{j = 1}^n g_j \partial_j (h)  \right)
}
{ =} { \sum_{i = 1}^n f_i \left(  \sum_{j = 1}^n \left(  \partial_i (g_j) \partial_j (h) +  g_j \partial_i \partial_j (h)  \right)  \right)
}
{ =} { \sum_{j = 1}^n \left(  \sum_{i = 1}^n f_i  \partial_i g_j  \right) \partial_j h +  \sum_{j = 1}^n \sum_{i = 1}^n    f_i g_j \partial_i \partial_j h
}
}
{}
{}{.}

}

An dieser expliziten Beschreibung sieht man

\mavergleichskettedisp
{\vergleichskette
{ \left(  \sum_{i = 1}^n f_i \partial_i  \right) \circ  \left(  \sum_{j = 1}^n g_j \partial_j   \right)
}
{ \neq} { \left(  \sum_{j = 1}^n g_j \partial_j   \right) \circ  \left(  \sum_{i = 1}^n f_i \partial_i  \right)
}
{ } {
}
{ } {
}
{ } {
}
}
{}{}{,}
da zwar der zweite Summand gleich ist, aber nicht der erste. Dies hat auch die folgende Konsequenz.
__NOEDITSECTION__

\inputfaktbeweis
{Offene Menge/R^n/Vektorfelder/Lie-Klammer/Vektorfeld/Fakt}
{Lemma}
{}
{

\faktsituation {Es sei

\mavergleichskette
{\vergleichskette
{ U
}
{ \subseteq }{ \R^n
}
{  }{
}
{    }{
}
{   }{
}
}
{}{}{}
\definitionsverweis {offen}{}{}
und seien $f_1 , \ldots , f_n, g_1 , \ldots , g_n$ zweifach
\definitionsverweis {stetig differenzierbare Funktionen}{}{}
auf $U$ mit den zugehörigen Differentialoperatoren

\mavergleichskette
{\vergleichskette
{ D
}
{ = }{ \sum_{i = 1}^n f_i \partial_i
}
{  }{
}
{    }{
}
{   }{
}
}
{}{}{}
und

\mavergleichskette
{\vergleichskette
{ E
}
{ = }{ \sum_{j = 1}^n g_j \partial_j
}
{  }{
}
{    }{
}
{   }{
}
}
{}{}{}}
\faktfolgerung {Dann gilt

\mavergleichskettedisp
{\vergleichskette
{ D \circ E -E \circ D
}
{ =} { \sum_{j = 1}^n \left(  \sum_{i = 1}^n \left(  f_i  \partial_i g_j - g_i \partial_i f_j  \right)  \right)  \partial_j
}
{ } {
}
{ } {
}
{ } {
}
}
{}{}{.}
Insbesondere entspricht dieser Ausdruck selbst einem Differentialoperator der Ordnung $1$ und somit einem Vektorfeld.}
\faktzusatz {}
\faktzusatz {}

}
{

Nach
[[Offene Menge/R^n/Vektorfelder/Hintereinanderschaltung/Explizit/Fakt|Lemma 11.5]]
ist

\mavergleichskettealignhandlinks
{\vergleichskettealignhandlinks
{D \circ E-E \circ D
}
{ =} { \left(  \sum_{i = 1}^n f_i \partial_i  \right) \circ  \left(  \sum_{j = 1}^n g_j \partial_j   \right) -  \left(  \sum_{j = 1}^n g_j \partial_j   \right) \circ  \left(  \sum_{i = 1}^n f_i \partial_i  \right)
}
{ =} { \sum_{j = 1}^n  \left(  \sum_{i = 1}^n  \left(  f_i  \partial_i g_j - g_i \partial_i f_j  \right)  \right)  \partial_j
}
{ } {
}
{ } {
}
}
{}
{}{.}

}

Auf einer Mannigfaltigkeit bezeichnen wir den Differentialoperator zu einem Vektorfeld $V$ mit $D_V$. Zu einer differenzierbaren Funktion
\maabb {f} { M } {\R
} {}
ist

\mavergleichskette
{\vergleichskette
{ D_V(f)
}
{ = }{ T(f) \circ V
}
{  }{
}
{    }{
}
{   }{
}
}
{}{}{.}
Angewendet auf einen Punkt

\mavergleichskette
{\vergleichskette
{ P
}
{ \in }{ M
}
{  }{
}
{    }{
}
{   }{
}
}
{}{}{}
ist

\mavergleichskettedisp
{\vergleichskette
{ \left(  D_V(f)  \right) (P)
}
{ =} { T_P(f) (V(P))
}
{ } {
}
{ } {
}
{ } {
}
}
{}{}{.}
Entsprechendes gilt für eine auf einer offenen Teilmenge von $M$ definierten differenzierbaren Funktion.

\inputdefinition
{}
{

Es seien
\mathkor {} {V} {und} {W} {}
stetig differenzierbare
\definitionsverweis {Vektorfelder}{}{}
auf einer
$C^2$-\definitionsverweis {differenzierbaren Mannigfaltigkeit}{}{}
mit den zugehörigen Differentialoperatoren $D_V$ und $D_W$. Dann nennt man das Vektorfeld, das der Hintereinanderschaltung
\mathl{D_V \circ D_W -D_W \circ D_V}{} entspricht, die
\definitionswort {Lie-Klammer}{}
der beiden Vektorfelder. Sie wird mit $[V,W]$ bezeichnet.

}

Es ist eine Konvention, welche Reihenfolge der Operatoren mit einem Minuszeichen in die Definition eingeht.
__NOEDITSECTION__

\inputfaktbeweis
{Mannigfaltigkeit/Vektorfelder/Lie-Klammer/Eigenschaften/Fakt}
{Lemma}
{}
{

\faktsituation {Es sei $M$ eine
$C^2$-\definitionsverweis {differenzierbare Mannigfaltigkeit}{}{.}}
\faktuebergang {Dann erfüllt die
\definitionsverweis {Lie-Klammer}{}{}
von
\definitionsverweis {Vektorfeldern}{}{}
die folgenden Eigenschaften.}
\faktfolgerung {\aufzaehlungdrei{Es ist

\mavergleichskette
{\vergleichskette
{ [V,W]
}
{ = }{ -[W,V]
}
{  }{
}
{    }{
}
{   }{
}
}
{}{}{.}
}{Es ist

\mavergleichskette
{\vergleichskette
{ [V_1+V_2,W]
}
{ = }{ [V_1,W] + [V_2,W]
}
{  }{
}
{    }{
}
{   }{
}
}
{}{}{.}
}{Für eine differenzierbare Funktion $f$ auf $M$ ist

\mavergleichskettedisp
{\vergleichskette
{ [fV,W]
}
{ =} { f [V,W] - \left(  D_Wf  \right)  V
}
{ } {
}
{ } {
}
{ } {
}
}
{}{}{.}
}}
\faktzusatz {}
\faktzusatz {}

}
{

(1) folgt aus der Definition, (2) folgt aus der Beziehung

\mavergleichskettedisp
{\vergleichskette
{ D_{V_1 +V_2 }
}
{ =} { D_{V_1 } + D_{V_2 }
}
{ } {
}
{ } {
}
{ } {
}
}
{}{}{}
zwischen Differentialoperatoren und Vektorfeldern. (3). Da die Aussage lokal ist, können wir

\mavergleichskettedisp
{\vergleichskette
{ V
}
{ =} { \sum_{i = 1}^n f_i \partial_i
}
{ } {
}
{ } {
}
{ } {
}
}
{}{}{}
und

\mavergleichskette
{\vergleichskette
{ W
}
{ = }{ \sum_{j = 1}^n g_j \partial_j
}
{  }{
}
{    }{
}
{   }{
}
}
{}{}{}
ansetzen. Dann ist

\mavergleichskettedisp
{\vergleichskette
{ fV
}
{ =} { f \sum_{i = 1}^n f_i \partial_i
}
{ =} { \sum_{i = 1}^n ff_i \partial_i
}
{ } {
}
{ } {
}
}
{}{}{}
und nach
[[Offene Menge/R^n/Vektorfelder/Lie-Klammer/Vektorfeld/Fakt|Lemma 11.6]]
ist

\mavergleichskettealignhandlinks
{\vergleichskettealignhandlinks
{ \, [fV,W]
}
{ =} { \sum_{j = 1}^n \left(  \sum_{i = 1}^n \left(  f f_i \partial_i g_j - g_i \partial_i (f f_j) \right)  \right) \partial_j
}
{ =} { \sum_{j = 1}^n \left(  \sum_{i = 1}^n \left(  f f_i \partial_i g_j - g_i \left(  f \partial_i ( f_j) +f_j \partial_i (f)  \right)  \right)  \right)  \partial_j
}
{ =} { f \sum_{j = 1}^n \left(  \sum_{i = 1}^n \left(  f_i \partial_i g_j - g_i \partial_i ( f_j)  \right)  \right) \partial_j - \sum_{j = 1}^n \left(  \sum_{i = 1}^n g_i f_j \partial_i (f)  \right)  \partial_j
}
{ =} { f[V,W] - \sum_{j = 1}^n \left(  \sum_{i = 1}^n g_i \partial_i (f)  \right) f_j \partial_j
}
}
{
\vergleichskettefortsetzungalign
{ =} { f[V,W] - \left(  \sum_{i = 1}^n g_i \partial_i (f)  \right) \left(  \sum_{j = 1}^n f_j \partial_j  \right)
}
{ =} { f[V,W] - D_W(f) V
}
{ } {}
{ } {}
}
{}{.}

}

 [[Kategorie:Latexseite]]
